\documentclass[10pt,a4paper]{amsart}
\usepackage[margin=19mm]{geometry}
\usepackage{amsmath,amssymb,mathtools}
\usepackage[scr=rsfs,cal=euler]{mathalfa}
\usepackage{xfrac}
\usepackage[unicode,hidelinks,bookmarks=false]{hyperref}
\newcommand{\ie}{i.e.\ }
\newcommand{\cf}{cf.\ }
\newcommand{\resp}{respectively\ }
\newcommand{\Subsection}{\subsection}
\providecommand{\nobreakdash}{\nobreak\hbox{-}}
\setlength{\emergencystretch}{2em}
\multlinegap0pt
\usepackage{bm}
\usepackage{bbm}
\usepackage{dsfont}
\usepackage{xspace}
\usepackage{cancel}
\usepackage{mathtools}
\usepackage{amsthm}
\makeatletter
\@ifundefined{thm}{\newtheorem{thm}{Thm}}{}
\@ifundefined{lem}{\newtheorem{lem}[thm]{Lemma}}{}
\@ifundefined{rem}{\newtheorem{rem}[thm]{Remark}}{}
\makeatother
\title[Matrices with simple symmetric digraphs and their group inve]{Matrices with simple symmetric digraphs and their group inverses}
\author{Raju Nandi}
\date{Source: arXiv:2312.08691v3 (CC BY 4.0)}
\begin{document}
\maketitle

\noindent\emph{Source and licence.} Matrices with simple symmetric digraphs and their group inverses. Version arXiv:2312.08691v3 (CC BY 4.0), distributed under the Creative Commons Attribution 4.0 International licence, as recorded in the arXiv licence metadata for this exact version and shown on the abstract page https://arxiv.org/abs/2312.08691v3. Full rights notice: licenses/arxiv-2312.08691-CC-BY-4.0.txt.\par\smallskip

\noindent\textit{Editorial scope.} Counted unit: the Lem whose source label is \texttt{1} in the article (source0.tex lines 430--432, source label \texttt{1}), delivered together with the article's own complete original proof of it (source0.tex lines 433--471, one \texttt{proof} environment of 39 lines). No mathematics is written, repaired or re-derived: statement and proof are the article's own text line for line. Required context retained uncounted: properties (lines 352--354). Every definition, display and label of those lines is retained unchanged. Omitted from this package and not redistributed: 1--351, 355--429, 472--802. Nothing inside a retained excerpt is omitted or altered: every retained body is delivered exactly as the article's lines are written, so no omission receipt of any kind accompanies this package. Citations: the retained excerpts carry no citation command at all, so no bibliography entry of the article is transcribed and no reference is redistributed. Reference adaptation: none was needed and none was made; every reference inside the delivered unit resolves inside it, so no unresolved reference or citation is produced. Printed slips kept literal: no printed word, symbol, hypothesis or number of the counted statement or of its proof was altered. Layout and class: the article's own class is \texttt{amsart} with packages including amsmath, amssymb, amsfonts, amstext, amsthm, color, tocvsec2, hyperref, euscript, enumerate, geometry, float, pgf, tikz; the wrapper is the lane's pinned \texttt{amsart} editorial wrapper at 10pt on A4, which restates the theorem-like environments the retained text needs and adds the article's own macro definitions that the retained text needs (none), with extra wrapper packages (bm, bbm, dsfont, xspace, cancel, mathtools). The delivered numbering is the unit's own. Assets: no retained span contains a figure environment, a graphics-inclusion command, a PDF-page inclusion, a table or a TikZ picture; no image file is copied into this package, the package declares no asset dependency and no image file is redistributed with it. Licence: version arXiv:2312.08691v3 of this article is distributed under the Creative Commons Attribution 4.0 International licence, as recorded in the arXiv licence metadata for this exact version and shown on the abstract page https://arxiv.org/abs/2312.08691v3. A full rights notice is delivered at licenses/arxiv-2312.08691-CC-BY-4.0.txt. Characters: every retained excerpt is pure ASCII, so the delivered engine, XeLaTeX with its default Latin Modern text font, reports no missing character.
\begin{rem}\label{properties}
Let $D\in \mathcal{D}$ have $k$ non-pendant vertices. Then, the number of 2-cycles in a maximum matching is always $k$. Note that, every non-pendant vertex is matched in any maximum matching of $D$, and both the endpoints of a length three alternating cycle chain are pendant vertices and a length one alternating cycle chain is nothing but a pendant cycle.
\end{rem}

\begin{lem}\label{1}
Let $A=(a_{ij})$ be an $n\times n$ real matrix such that $D(A)\in \mathcal{D}$. Let $D(A)$ be a strongly connected digraph and $\Delta _A \neq 0$. Let $B=(b_{ij})$ be the matrix given by $b_{ij}=\frac{\mu _{ij}}{\Delta _A}, 1\leq i,j\leq n$. Then, $AB=BA$.
\end{lem}

\begin{proof}
Let $A=(a_{ij})$. Then, $AB=BA$ has the following equivalent form :
\begin{equation}\label{ab=ba}
\sum _{k=1}^na_{ik}\mu _{kj}=\sum _{l=1}^n\mu _{il} a_{lj}\  \  \  for~every\  \  i,j\in \{1,2,\ldots n\}.
\end{equation}
First, we discuss the case $i=j$. Let $\{i_1,i_2,\ldots ,i_t\}\in \{1,2,\ldots ,n\}$ be such that for any $m\in \{1,2,\ldots ,t\}$, $a_{i,i_m}\neq 0$ and the cycle $(i,i_m,i)$ belongs to some maximum matching in $D(A)$. Since $\mathbb{M}(i,i_m)=\mathbb{M}(i_m,i)$ and $\beta _{\overline{i,i_m}}(M)=\beta _{\overline{i_m,i}}(M)$, the expressions on both the sides of equation (\ref{ab=ba}) are equal and they equal the common value $\sum_{m=1}^t\big (a_{ii_m}a_{i_mi}\sum _{M\in \mathbb{M}(i_m,i)}\beta _{\overline{i_m,i}}(M)\big)=\sum _{M\in \mathbb{M}(i)}\eta (M)$.

Now, we consider the fact $i\neq j$ and prove \ref{ab=ba} by considering four cases.\\
{\bf Case (i)}: $i$ and $j$ are pendant vertices.\\ 
\underline{{\it Subcase \rm{(i.1)}}}: $i$ and $j$ have a unique common neighbour. Let $q$ be the common neighbour. Then, $(i,q,i)$ and $(j,q,j)$ are both maximum matching cycles and cannot simultaneously be present in a maximum
matching. So, $\{M\backslash \{(i,q,i)\}\vert ~M\in \mathbb{M}(i,q)\}=\{M\backslash \{(j,q,j)\}\vert ~M\in \mathbb{M}(j,q)\}$. Thus,
\begin{align*}
a_{iq}\mu _{qj}& =a_{iq}a_{qj}\sum _{M\in \mathbb{M}(q,j)}\beta _{\overline{q,j}}(M) \\
 &= a_{qj}\Big (a_{iq}\sum _{M\in \mathbb{M}(i,q)}\beta _{\overline{i,q}}(M)\Big ) \\
 &= \mu _{iq}a_{qj}. 
\end{align*} 
It follows that equation \ref{ab=ba} holds.\\
\underline{{\it Subcase \rm{(i.2)}}}: There is no common neighbour for the vertices $i$ and $j$. Let vertices $i$ and $j$ be adjacent to non-pendant vertices $q$ and $p$, respectively. Then, the length of any cycle chain from $i$ to $j$ is at least three. As we have observed in Remark \ref{properties}, both the end vertices of a length three alternating cycle chain must be pendant vertices, $\mu _{qj}$ and $\mu _{ip}$ are both zero. So, equation \ref{ab=ba} is vacuously true.\\
{\bf Case (ii)}: $i$ and $j$ are non-pendant vertices. Since a non-pendant cycle can not be present in a maximum matching and the end vertices of a length three alternating cycle chain are pendant, $\mu _{qj}$ and $\mu _{ip}$ are zero for any arbitrary vertex $q$ adjacent to $i$ and $p$ adjacent to $j$ respectively. So, the expressions on both sides of equation \ref{ab=ba} are zero.\\
{\bf Case (iii)}: $i$ is a pendant vertex, while $j$ is a non-pendant vertex. Let $N_j=\{j_1,j_2,\ldots ,j_s\}$ be the set of all pendant vertices adjacent to $j$. We consider two subcases.\\
\underline{{\it Subcase \rm{(iii.1)}}}: $i$ is adjacent to $j$. then, the left hand side of \ref{ab=ba} is zero and the right hand side equal to
\begin{equation}\label{split}
\sum _{l=1}^s\mu _{ij_l} a_{j_lj} + \sum _{\{m\vert m \in N(j)\backslash N_j\}}\mu _{im} a_{mj}.
\end{equation}
% $$\sum _{l=1}^s\mu _{ij_l} a_{j_lj} + \sum _{\{m\vert m \in N(j)\backslash N_j\}}\mu _{im} a_{mj}.$$
The first sum is zero because $i$ and $j_l$ are pendant vertices with common neighbour $j$ for all $l\in \{1,2,\ldots ,s\}$. In the second sum, for any $m$, the vertex $m$ is non-pendant and also it is not adjacent to pendant vertex $i$. So, $\mu _{im}$ is zero.\\
\underline{{\it Subcase \rm{(iii.2)}}}: $i$ is not adjacent to $j$. Let $q$ be the non-pendant vertex adjacent to $i$. The analysis in this case is split into two further subcases, which we consider next. \\
\underline{{\it Subcase \rm{(iii.2.a)}}}: $q$ is adjacent to $j$. Since non-pendant cycles do not belong to any maximum matching, again, the left hand side of \ref{ab=ba} is zero, and the right hand side can be split as \ref{split}. Let $\mathbb{M}_{j_l}(i,q)\subseteq \mathbb{M}(i,q)$ be the set of maximum matching containing pendant cycle $(j,j_l,j)$. Then $\mathbb{M}(i,q)=\cup _{l=1}^s\mathbb{M}_{j_l}(i,q)$, a mutually disjoint union. Now, the second sum of \ref{split} is
\begin{align*}
\mu _{iq}a_{qj}& =a_{iq}a_{qj}\sum _{M\in \mathbb{M}(i,q)}\beta _{\overline{i,q}}(M) \\
 &= a_{qj}a_{iq}\sum _{l=1}^s\sum _{M\in \mathbb{M}_{j_l}(i,q)}\beta _{\overline{i,q}}(M) \\
 &= a_{qj}a_{iq}\sum _{l=1}^sa_{jj_l}a_{j_lj}\sum _{M\in \mathbb{M}_{j_l}(i,q)}\beta _{\overline{i,j_l}}(M) \\
 &= -\sum _{l=1}^sa_{j_lj}\big (-a_{iq}a_{qj}a_{jj_l}\big )\sum _{M\in \mathbb{M}(i,j_l)}\beta _{\overline{i,j_l}}(M) \\
 &= -\sum _{l=1}^sa_{j_lj}\mu _{ij_l}.
\end{align*} 
So, the right hand side of \ref{ab=ba} is also zero.\\
\underline{{\it Subcase \rm{(iii.2.b)}}}: $q$ is not adjacent to $j$. Since $q$ and $j$ both are non-pendant vertices, the left hand side of equation \ref{ab=ba} is zero. After splitting the right hand side as \ref{split}, in the first sum, the length of any cycle chain from vertex $i$ to $j_l$ must be at least four for any $l\in \{1,2,\ldots ,s\}$ and so, $\mu _{ij_j}$ is zero. Next, as $m$ is a non-pendant vertex and since it is not adjacent to $i$, the second sum is also zero, showing that the right hand side equals zero.\\
{\bf Case (iv)}: $i$ is a non-pendant vertex and $j$ is a pendant vertex. Let $N_i=\{i_1,i_2,\ldots ,i_r\}$ be the set of all pendant vertices adjacent to $i$. Then, the proof is the same as in Case (iii) by interchanging the roles of $i$ and $j$.
\end{proof}

\end{document}
